
\usepackage{mathtools}
\usepackage{graphicx}
\usepackage{subcaption}
\usepackage{caption}
\usepackage{booktabs}
\usepackage{adjustbox}
\usepackage{multirow}
\usepackage{caption}
\usepackage{eurosym} 
\usepackage{xcolor}
\usepackage{comment}
\usepackage{subcaption}
\usepackage{amssymb, amsmath}
\usepackage{algorithm}
\usepackage{algpseudocode}
\usepackage{blindtext}
\usepackage{pdflscape}
\usepackage{rotating}
\usepackage{longtable}
\usepackage{placeins}
\usepackage{spverbatim}
\usepackage{tikz}
\usepackage{rotating}
\usepackage{pgfplots}
\pgfplotsset{width=10cm,compat=1.9}

\usepackage{makecell}
\usepackage[flushleft]{threeparttable}
\renewcommand\theadalign{bc}
\renewcommand\theadfont{\bfseries}
\renewcommand\theadgape{\Gape[4pt]}
\renewcommand\cellgape{\Gape[4pt]}
\usepackage{tikz}
\usepackage{tikz-3dplot}

\usetikzlibrary{shapes,shapes.geometric, positioning, fit, backgrounds, calc, 3d, shapes.multipart, arrows.meta, patterns, perspective}

\newcommand\lbo{0.5}
\newcommand\wbo{1}
\newcommand\hbo{1}
\newcommand\lbt{2}
\newcommand\wbt{1}
\newcommand\hbt{1}
\newcommand\separationofcircles{0.1}


\tikzset{
    single_redbox/.pic={
    \draw[red, dashed, opacity=0.5] (0-\lbo,0-\wbo,0-\hbo) -- ++(0,0,\hbo);
    \draw[red, dashed, opacity=0.5] (0-\lbo,0-\wbo,0-\hbo) -- ++(0,\wbo,0);
    \draw[red, dashed, opacity=0.5] (0-\lbo,0-\wbo,0-\hbo) -- ++(\lbo,0,0);
    
    % Draw front faces
    \draw[red, fill=red!20, opacity=0.5] (0,0,0) -- ++(-\lbo,0,0) -- ++(0,-\wbo,0) -- ++(\lbo,0,0) -- cycle;
    \draw[red, fill=red!30, opacity=0.5] (0,0,0) -- ++(0,0,-\hbo) -- ++(0,-\wbo,0) -- ++(0,0,\hbo) -- cycle;
    \draw[red, fill=red!40, opacity=0.5] (0,0,0) -- ++(-\lbo,0,0) -- ++(0,0,-\hbo) -- ++(\lbo,0,0) -- cycle;
    }
}
\tikzstyle{line} = [draw, -Latex]

\newcommand{\neworrenewcommand}[1]{\providecommand{#1}{}\renewcommand{#1}}

\newcommand{\drawblock}[9]{
    \neworrenewcommand{\ffoo}[1]{
           
        \foreach \i in {0,...,#8} {
            \foreach \j in {0,...,#9} {
                \foreach \k in {0,...,##1} {
                    % Only draw the visible faces
                    % Front face
                    \ifnum \i=#8
                        \draw[#7, fill=#7!50] (#1+\i*#4+#4,#2+\j*#5,#3+\k*#6) -- ++(0,#5,0) -- ++(0,0,#6) -- ++(0,-#5,0) -- cycle;
                    \fi
                    % Right face
                    \ifnum \j=#9
                        \draw[#7, fill=#7!30] (#1+\i*#4,#2+#5+\j*#5,#3+\k*#6) -- ++(#4,0,0) -- ++(0,0,#6) -- ++(-#4,0,0) -- cycle;
                    \fi
                    % Top face
                    \ifnum \k=##1
                        \draw[#7, fill=#7!20] (#1+\i*#4,#2+\j*#5,#3+#6+\k*#6) -- ++(#4,0,0) -- ++(0,#5,0) -- ++(-#4,0,0) -- cycle;
                    \fi
                    % Draw a point at each vertex
                    % \fill (#1+\i*#4,#2+\j*#5,#3+\k*#6) circle (1pt);
                }
            }
        }
    
    }
    \ffoo
}



      \tikzset{
    box/.style={rectangle, draw, fill=gray!20, align=center, minimum height=1cm, minimum width=2.5cm},
    ellip/.style={draw, ellipse, align=center, minimum width=2.5cm, minimum height=1cm},
    label_node/.style={circle, draw, fill=#1, text=black, inner sep=1pt, minimum size=1em, font=\bfseries},
    gene_node/.style={ align=center, inner sep=2pt, fill=white, minimum width=0.5cm, minimum height=1cm},
    gene_group/.style={rectangle, draw, inner sep=6pt, rounded corners, align=center, minimum height=2.5cm},
    aux_gene_group/.style={align=center, inner sep=2pt, minimum width=2cm, minimum height=1cm}
  }

\providecommand{\keywords}[1]
{
  \small	
  \textbf{\textit{Keywords---}} #1
}

% Visible marker for a number or table that does not exist yet.  Every use is
% a place where a generated fragment or a regenerated run must supply the
% value; grep for \pending before submitting.
\providecommand{\pending}[1]{\textcolor{red}{\textbf{[pending: #1]}}}
\providecommand{\pendingcell}{\textcolor{red}{\textbf{[pending]}}}
